\section{Charges, completeness and enclosing area}\label{sec:transfer}

We now finish Theorem~\ref{prop:replacement}. All estimates in this
section concern the data constructed in Sections~\ref{sec:collar}
and~\ref{sec:bending}; in particular their DEC has already been proved.
Subscripts $R$ are suppressed when they are not needed.

\subsection{The asymptotically flat output}
For each fixed sufficiently large $R$, the outer construction has
$v=\beta=0$ beyond a finite radius. Therefore $L=0$, $\tau=0$ and
$k=0$ there, and
\begin{equation}\label{eq:af-tail}
 G=(1-2F/r)^{-1}\mathrm dr^2+r^2\sigma.
\end{equation}
Lemma~\ref{lem:bending} bounds every fixed-order scaled radial and
angular derivative of $F$. In Cartesian coordinates $x=rn$ the metric
takes the form
\[
 G_{ij}=\delta_{ij}+b\,n_i n_j,
 \qquad b=(1-2F/r)^{-1}-1.
\]
Every Cartesian derivative can be expressed as
$\partial_{x_i}=r^{-1}(n_iD+X_i)$, where $D=r\partial_r$ and
$X_i$ is a smooth tangential field on the unit sphere. Repeatedly applying
this identity to $b n_i n_j$ proves
$G_{ij}-\delta_{ij}=O_j(r^{-1})$ for every fixed finite $j$.
In particular the $O_6$ metric requirement of the numerical input below
holds. The round constraint estimates
give $C=O(r^{-7/2})$ and $J=0$ on this tail. Since its volume element
is comparable to $r^2\mathrm dr\,\mathrm dA_\sigma$, the constraint
densities are integrable. The portion inside a fixed tail sphere is
compact and smooth, so integrability holds on the whole exterior.

For completeness, the ADM energy calculation allows angular dependence
of $b$. Direct differentiation, using $\partial_jn_i=(\delta_{ij}
-n_i n_j)/r$, gives
\[
 (\partial_jG_{ij}-\partial_iG_{jj})n^i=\frac{2b}{r}.
\]
The derivatives of $b$ cancel after contraction with $n$; there is no
assumption that $F$ is radial. Since $rb=2F+O(r^{-1})$, uniformly on
the sphere, \eqref{eq:adm-energy} gives
\begin{equation}\label{eq:replacement-charges}
 E_R=\lim_{r\to\infty}\avg{F(r,\cdot)}
 =m_B+c_R(\infty),\qquad P_R=0.
\end{equation}
The mean of the heat solution is constant, which proves the second
equality. The vanishing of $k$ on the tail proves $P_R=0$ directly.
The bound in \eqref{eq:mass-padding} yields
$0\leq c_R(\infty)\leq2R^{-1/2}$. Because $m_B>0$, the constructed
charges satisfy the strict inequality $E_R>|P_R|$.
These are the charges of one completed replacement, with $R$ fixed.
The constant in the refined tail estimate~\eqref{eq:tail-profile}
need not be uniform in $R$.

\subsection{A comparison for every cut}
In the balanced chart the original metric satisfies
$|g-h|_h\leq C_* r^{-3}$ for $r\geq R$. On $R\leq r\leq3R$,
Lemma~\ref{lem:collar} gives the same estimate for $G$. Beyond $3R$,
the new angular metric is exactly $r^2\sigma$ and has no shift, while
\[
 u^2=1+v^2-2F/r\leq1+r^2+2\|F\|_\infty/r.
\]
Write $M_*\geq\|F\|_\infty$ for a bound independent of $R$.
Since $0\leq v\leq r$, comparison of the radial coefficients gives
\[
 u^{-2}\geq\frac{(1+r^2)^{-1}}{1+2M_*r^{-3}}
       \geq(1-2M_*r^{-3})(1+r^2)^{-1}.
\]
The angular coefficients already agree with those of $h$. Thus
$G\geq(1-2M_*r^{-3})h$ on the whole outer region, while
$g\leq(1+C_*r^{-3})h$. Together with the collar comparison and
the identity $G=g$ on the compact interior, this gives a constant
$C_{\mathrm{met}}$ independent of $R$ such that, with
$\varepsilon_R=C_{\mathrm{met}}R^{-3}$,
\begin{equation}\label{eq:metric-comparison}
 G_R\geq(1-\varepsilon_R)g\quad\hbox{on all of }\Omega.
\end{equation}
We increase $R$ once so that $0\leq\varepsilon_R<1$.

This also proves completeness, including the boundary. A $G_R$-Cauchy
sequence is $g$-Cauchy by \eqref{eq:metric-comparison}, and hence has a
limit in $\Omega$. The smooth positive metrics induce the same local
topology, so it converges in $G_R$ as well. The manifold, its boundary,
and its single end have not changed. Compact truncations remain compact
after the balancing chart change because old and new radii are
comparable sufficiently far out.

On the tangent plane of any smooth cut, the two relative eigenvalues
of $G_R$ with respect to $g$ are at least $1-\varepsilon_R$. Its area
element is therefore at least $(1-\varepsilon_R)$ times the original
area element. The admissible cuts are the same geometric surfaces for
the two metrics, with the same filled-side convention. Taking the
infimum over \emph{all} such surfaces gives
\begin{equation}\label{eq:area-comparison}
 A_{\min,G_R}(S)\geq(1-\varepsilon_R)A_{\min,g}(S).
\end{equation}
This includes cuts coinciding with the obstacle $S$. It uses neither
existence of a minimizing cut nor any identification of its area with
the area of $S$. Equations~\eqref{eq:replacement-charges}--
\eqref{eq:area-comparison}, together with the smooth DEC construction,
prove Theorem~\ref{prop:replacement}.

\subsection{The numerical input and the Bondi bound}\label{sec:companion-input}
The replacements now have the geometric properties needed for the
numerical step. We use the dimension-three case of the companion
\cite[Definitions~1.1--1.2 and Theorem~1.3]{NumericalPenrose}, stated
here with our physical constraint normalization.

\begin{theorem}[Asymptotically flat numerical input]
\label{ass:af}
Let $(X,q,k)$ be a smooth connected orientable three-dimensional
exterior, complete including its nonempty compact smooth boundary,
with exactly one asymptotically flat coordinate end and compact
complement. Suppose that, in Cartesian coordinates, for some
$1/2<a<1$,
\[
 q_{ij}-\delta_{ij}=O_6(r^{-a}),\qquad
 k_{ij}=O_5(r^{-1-a}).
\]
Use the constraints \eqref{eq:constraints-definition} and ADM
normalizations \eqref{eq:adm-energy}--\eqref{eq:adm-momentum}.
Assume DEC, integrability of $\mu$ and $|J|_q$, finite ADM limits,
$E>|P|$, and $A_{\min,q}(\partial X)>0$. If every boundary component
satisfies $\theta_+\leq0$ with normal into $X$, then
\[
 \sqrt{E^2-|P|^2}\geq
 \sqrt{\frac{A_{\min,q}(\partial X)}{16\pi}}.
\]
The full cut class is that of \eqref{eq:area-definition}, including
disconnected cuts and partial or whole coincidence with the boundary.
No extension across the boundary, outermostness, outer
area-minimization or attainment of the infimum is assumed.
\end{theorem}

To identify this statement with the cited theorem, its densities satisfy
\[
 2\mu_{\mathrm N}=R+(\tr k)^2-|k|^2,\qquad
 J_{\mathrm N}=\diver(k-(\tr k)q).
\]
Thus $\mu_{\mathrm N}=8\pi\mu$ and $J_{\mathrm N}=8\pi J$;
the common positive factor preserves DEC and integrability. At
dimension three its ADM factors are
$[2(n-1)\omega_{n-1}]^{-1}=1/(16\pi)$ and
$[(n-1)\omega_{n-1}]^{-1}=1/(8\pi)$, so the charges are unchanged.
Its Definition~1.2 takes the full intrinsic boundary of the connected
closed end-containing outer domain, including every component and
all contact with the original boundary. This is precisely the cut
class in \eqref{eq:area-definition}.

\begin{proof}[Proof of Corollary~\ref{cor:cks-inequality}]
Fix a sufficiently large $R$. The data $(G_R,k_R)$ have the same
connected orientable manifold, compact boundary and single end as the
original data. Their completeness follows from
\eqref{eq:metric-comparison}. Both tensors agree with $(g,K)$ on a
neighborhood of $S$, so its normal, mean curvature, tangential tensor
trace and weak future trapping are unchanged. The construction proves
DEC and integrable constraint densities. On the completed tail,
$G_R-\delta=O_6(r^{-1})$ and $k_R=0$, which give the required
differentiated rates with $a=3/4$. Equation~\eqref{eq:area-comparison}
and the assumed positivity of $A_{\min,g}(S)$ give
$A_{\min,G_R}(S)>0$. Finally,
\eqref{eq:replacement-charges} gives finite charges with
$E_R=m_B+c_R(\infty)>0=|P_R|$.

Thus every hypothesis of the companion's numerical theorem, in the
form of Theorem~\ref{ass:af}, holds for this fixed replacement.
Using the comparison of all cuts gives
\begin{equation}\label{eq:finite-numerical}
 m_B+c_R(\infty)
 \geq\sqrt{\frac{A_{\min,G_R}(S)}{16\pi}}
 \geq\sqrt{\frac{(1-\varepsilon_R)A_{\min,g}(S)}{16\pi}}.
\end{equation}
Only now let $R\to\infty$, using $0\leq c_R(\infty)\leq2R^{-1/2}$
and $\varepsilon_R\to0$. This proves \eqref{eq:main}. The examples
in Section~\ref{sec:sharpness} establish sharpness.
\end{proof}

There are two successive infinities in this proof. Radial infinity
defines the ADM charges of one completed replacement. The replacement
radius tends to infinity only after its numerical inequality has
been proved. No convergence of comparison metrics or minimizing cuts
is required, and the constant in $O_R(r^{-3})$ need not be uniform in
the replacement radius.

\begin{remark}[The weaker-decay numerical theorem]
The companion also proves the same inequality directly under
$q-\delta=O_2(r^{-a})$, $k=O_1(r^{-1-a})$, $a>1/2$, for the same
smooth one-ended exterior and full cut class, with physical constraint
densities, integrable DEC, finite timelike charges, and weakly future
trapped boundary \cite[Definition~7.1 and Theorem~7.2]{NumericalPenrose}.
That theorem does not assume positivity of the enclosing infimum.
It gives another numerical input for our replacements; its weaker
original-data hypotheses are a separate result, not a consequence of
changing the exponent in Theorem~\ref{ass:af}.
\end{remark}

\paragraph{Dependence on the numerical input.}\label{sec:conditional}
The proof uses the companion only through the displayed asymptotically
flat inequality. More generally, this inequality is the sole numerical
input needed for the same argument to give the Bondi bound for every weakly future trapped,
possibly disconnected boundary under the other hypotheses of
Corollary~\ref{cor:cks-inequality}. In particular, it applies when $S$
is connected and marginal, and the open exterior contains no compact
smooth embedded two-sided separating union of closed future or past
apparent horizons, allowing either sign on each component. These
additional horizon conditions are not needed for the transfer.
The Schwarzschild examples below satisfy them and attain the bound
at every positive mass.

The inequality alone does not classify equality for the original data:
the construction need not preserve an equality case at finite $R$.
A classification, or an extension to weaker null-infinity expansions,
would require a further argument.
